% !TeX root = ../hodge-classes.tex
\begin{thebibliography}{99}
\bibitem{Andre1992}
Y. Andr\'e,
\emph{Une remarque \`a propos des cycles de Hodge de type CM},
in S\'eminaire de Th\'eorie des Nombres, Paris 1989--1990, Progress in Mathematics 102, Birkh\"auser, 1992, 1--7.

\bibitem{daSilva2021Fermat}
G.~da Silva Jr., \emph{Notes on the Hodge conjecture for Fermat varieties},
Experimental Results \textbf{2} (2021), e22.
\href{https://doi.org/10.1017/exp.2021.14}{doi:10.1017/exp.2021.14}.
Published Theorem~3.3, its preceding full-$P_{21}$ assertion, and Appendix.

% Add to the existing bibliography only when integrating mt_addition.tex.
% Official author-maintained source; Theorem 6.8 and Corollary 6.9,
% printed p.73, give the integral/rational covariant H_1 equivalences.

\bibitem{Markman2025}
E. Markman,
\emph{Cycles on abelian $2n$-folds of Weil type from secant sheaves on abelian $n$-folds},
arXiv:2502.03415v2 (2025), Section~1.1, p.~3, and Theorem~1.5.1, p.~9.
The theorem proves algebraicity of Hodge--Weil classes for polarized
abelian sixfolds of Weil type with discriminant $-1$, without a
principal-polarization hypothesis.
\url{https://arxiv.org/pdf/2502.03415v2}

\bibitem{MilneShimura2017}
J.~S.~Milne, \emph{Introduction to Shimura Varieties},
revised September 16, 2017, Theorem~6.8 and Corollary~6.9,
\url{https://www.jmilne.org/math/xnotes/svi.pdf}.

\bibitem{MilneCM2020}
J. S. Milne,
\emph{Complex Multiplication}, version 0.10, 14 July 2020,
Section 3.11 and Propositions 3.12--3.13 and 3.17, pp.~29--31.
\url{https://www.jmilne.org/math/CourseNotes/CM.pdf}

\bibitem{MilneHodge2020}
J. S. Milne,
\emph{Hodge classes on abelian varieties},
arXiv:2010.08857, 2020.
\url{https://arxiv.org/abs/2010.08857}

\bibitem{OpenAICM2026}
OpenAI math repository,
\emph{The rational Hodge conjecture for CM abelian varieties},
public manuscript dated 6 October 2026, introduction (main CM theorem)
and tensor-reduction section (auxiliary-factor contractions).
The comparison uses repository commit
\texttt{4eaee87548250b2682affda7d05d947c05554bae};
the general theorem and its formalization are not proof inputs here.
\url{https://github.com/openai/math/tree/4eaee87548250b2682affda7d05d947c05554bae/preprints/The-rational-Hodge-conjecture-for-CM-abelian-varieties-October-6-2026}

\bibitem{Saller2026}
K. Saller,
\emph{Split Weil structures on abelian sixfolds and the Hodge conjecture},
version 2, Zenodo record 23290280 (2026), Remarks~4.7--4.8.
\url{https://doi.org/10.5281/zenodo.23290280}

\bibitem{Schoen1998}
C. Schoen,
\emph{Addendum to ``Hodge classes on self-products of a variety with an automorphism''},
Compositio Math. 114 (1998), no.~3, 329--336.

\bibitem{Shioda1981}
T.~Shioda, \emph{Algebraic Cycles on Abelian Varieties of Fermat Type},
Mathematische Annalen \textbf{258} (1981), 65--80.
\href{https://doi.org/10.1007/BF01450347}{doi:10.1007/BF01450347}.
The original numbering used here is Lemmas~4.1--4.2, Theorems~4.3--4.4 and
Example~6.3; the latter concerns $J(y^2=x^{21}-1)$.
\url{https://gdz.sub.uni-goettingen.de/download/pdf/PPN235181684_0258/LOG_0017.pdf}

\end{thebibliography}
